\documentclass[10pt,a4paper]{amsart}
\usepackage[margin=19mm]{geometry}
\usepackage{amsmath,amssymb,mathtools}
\usepackage[scr=rsfs,cal=euler]{mathalfa}
\usepackage{xfrac}
\usepackage[unicode,hidelinks,bookmarks=false]{hyperref}
\newcommand{\ie}{i.e.\ }
\newcommand{\cf}{cf.\ }
\newcommand{\resp}{respectively\ }
\newcommand{\Subsection}{\subsection}
\providecommand{\nobreakdash}{\nobreak\hbox{-}}
\setlength{\emergencystretch}{2em}
\multlinegap0pt
\usepackage{bm}
\usepackage{bbm}
\usepackage{dsfont}
\usepackage{xspace}
\usepackage{cancel}
\usepackage{mathtools}
\usepackage{amsthm}
\makeatletter
\@ifundefined{theorem}{\newtheorem{theorem}{Theorem}}{}
\@ifundefined{assumption}{\newtheorem{assumption}{Assumption}}{}
\@ifundefined{lemma}{\newtheorem{lemma}[theorem]{Lemma}}{}
\makeatother
\renewcommand{\d}{{\rm{d}}}
\title[$(H,H^3)$-smoothing effect and convergence of solutions of s]{$(H,H^3)$-smoothing effect and convergence of solutions of stochastic two-dimensional anisotropic Navier-Stokes equations driven by colored noise}
\author{Hui Liu, Dong Su, Chengfeng Sun, Jie Xin}
\date{Source: arXiv:2508.16938v1 (CC BY 4.0)}
\begin{document}
\maketitle

\noindent\emph{Source and licence.} $(H,H^3)$-smoothing effect and convergence of solutions of stochastic two-dimensional anisotropic Navier-Stokes equations driven by colored noise. Version arXiv:2508.16938v1 (CC BY 4.0), distributed under the Creative Commons Attribution 4.0 International licence, as recorded in the arXiv licence metadata for this exact version and shown on the abstract page https://arxiv.org/abs/2508.16938v1. Full rights notice: licenses/arxiv-2508.16938-CC-BY-4.0.txt.\par\smallskip

\noindent\textit{Editorial scope.} Counted unit: the Lemma whose source label is \texttt{lem:H2bd} in the article (source0.tex lines 1150--1159, source label \texttt{lem:H2bd}), delivered together with the article's own complete original proof of it (source0.tex lines 1160--1378, one \texttt{proof} environment of 219 lines). No mathematics is written, repaired or re-derived: statement and proof are the article's own text line for line. Required context retained uncounted: assum (lines 167--173); 2.1 (lines 503--505); 3.9 (lines 596--598); 2.2 (lines 600--608); lemma4.3 (lines 947--955); 4.8 (lines 958--960); mar8.3 (lines 1123--1131). Every definition, display and label of those lines is retained unchanged. Omitted from this package and not redistributed: 1--166, 174--502, 506--595, 599--599, 609--946, 956--957, 961--1122, 1132--1149, 1379--3685. Nothing inside a retained excerpt is omitted or altered: every retained body is delivered exactly as the article's lines are written, so no omission receipt of any kind accompanies this package. Citations: the retained excerpts carry no citation command at all, so no bibliography entry of the article is transcribed and no reference is redistributed. Reference adaptation: none was needed and none was made; every reference inside the delivered unit resolves inside it, so no unresolved reference or citation is produced. Printed slips kept literal: no printed word, symbol, hypothesis or number of the counted statement or of its proof was altered. Layout and class: the article's own class is \texttt{elsarticle} with packages including lineno, times, hyperref, amsmath, amsthm, amscd, amsfonts, amssymb, graphicx, color, latexsym, geometry, mathabx; the wrapper is the lane's pinned \texttt{amsart} editorial wrapper at 10pt on A4, which restates the theorem-like environments the retained text needs and adds the article's own macro definitions that the retained text needs (\texttt{\textbackslash d}), with extra wrapper packages (bm, bbm, dsfont, xspace, cancel, mathtools). The delivered numbering is the unit's own. Assets: no retained span contains a figure environment, a graphics-inclusion command, a PDF-page inclusion, a table or a TikZ picture; no image file is copied into this package, the package declares no asset dependency and no image file is redistributed with it. Licence: version arXiv:2508.16938v1 of this article is distributed under the Creative Commons Attribution 4.0 International licence, as recorded in the arXiv licence metadata for this exact version and shown on the abstract page https://arxiv.org/abs/2508.16938v1. A full rights notice is delivered at licenses/arxiv-2508.16938-CC-BY-4.0.txt. Characters: every retained excerpt is pure ASCII, so the delivered engine, XeLaTeX with its default Latin Modern text font, reports no missing character.
\begin{assumption} \label{assum}
Let    $ h \in H^3$, and $\nabla\cdot h=0$ and
 \begin{align*}
  \|\nabla h\|_{L^\infty} <  \sqrt {\pi\delta} \frac{1}{2}\nu \lambda_1 ,
 \end{align*}
where  $\lambda_1 = 4\pi^2 /L^2$ is the first eigenvalue of the Stokes operator $A$.
  \end{assumption}

 \begin{align}\label{2.1}
   \partial_t\textbf{u}+ \nu A_1\textbf{u}+B(\textbf{u})  =f,\quad u(0)=u_{0}.
\end{align}

\begin{align}\label{3.9}
\|\nabla \textbf{v}\|^2=\|\nabla\times \textbf{v}\|^2\leq 2\|(\partial_yv_1,\partial_xv_2)\|^2.
\end{align}

\begin{align}\label{2.2}
\left\{
\begin{aligned}
&
\frac{\d \textbf{v}}{\d t}+ \nu A_1\textbf{v}+B\big( \textbf{v}(t)+hy_\delta(\theta_t\omega)\big)=f(\textbf{x})-  \nu A_1hy_\delta(\theta_t\omega)+hy_\delta (\theta_t\omega), \\
& \textbf{v}(0)=\textbf{u}_h(0)-hy_\delta (\omega).
\end{aligned}
\right.
\end{align}

\begin{lemma}[$H^1$-distance]
\label{lemma4.3}
 Let Assumption \ref{assum} hold and $f\in H$. There   exist   random variables $T_{\mathfrak B}(\cdot)$ and    $\zeta_4(\cdot)$, where  $\zeta_4(\cdot)$ is tempered,  such that the solutions $\textbf{v}$  of the random  anisotropic NS equations \eqref{2.2} driven by colored noise and $\textbf{u}$  of the deterministic anisotropic NS equations  \eqref{2.1} satisfy
\begin{align}
 \big \|A^{\frac 12} \textbf{v}(t,\theta_{-t}\omega,\textbf{v}(0))-A^{\frac 12} \textbf{u}(t,\textbf{u}(0)) \big\| ^2 \leq \zeta_4(\omega),\quad \forall t\geq T_{\mathfrak B}(\omega), \nonumber
\end{align}
for any $\textbf{v}(0)\in\mathfrak{B}(\theta_{-t}\omega)$
 and $\textbf{u}(0)\in\mathcal{A}_0$ for $\omega\in \Omega$.
\end{lemma}

\begin{align}\label{4.8}
\frac{\d w}{\d t}+ \nu  A_1w+B( \textbf{v}+hz(\theta_t\omega))-B(\textbf{u}) = hy_\delta(\theta_t\omega) - \nu A_1hy_\delta(\theta_t\omega).
\end{align}

\begin{align}
 &\sup_{\eta\in [t-2,t]}
  \left(
 \|A^{\frac12}\textbf{v} (\eta ,\theta_{-t } \omega, \textbf{v}(0)) \|^2
 + \frac{\nu}{2} \int_{\eta-2}^\eta  \|A\textbf{v} (s,\theta_{-t} \omega, \textbf{v}(0)) \|^2 \, \d s\right) \notag  \\
& \quad \leq  C\left (  \sup_{\tau\in (-5,0)}|y_\delta(\theta_\tau\omega)|^6+1\right)\big(1
 +   \zeta_1(\omega)   \big)
 =:  \zeta_3(\omega) ,\quad t\geq T_{\mathfrak B} (\omega) ,\label{mar8.3}
\end{align}

\begin{lemma}[$H^2$-distance] \label{lem:H2bd}  Let Assumption \ref{assum} hold and $f\in H$.
There exists a tempered random variable $\zeta_5(\cdot)$ such that the solutions $\textbf{v}$ of the random anisotropic NS equations \eqref{2.2} and $\textbf{u}$  of the deterministic anisotropic NS equations  \eqref{2.1} satisfy
\[
 \|A\textbf{v}(t,\theta_{-t}\omega,\textbf{v}(0))-A\textbf{u}(t,\textbf{u}(0))\|^2\leq \zeta_5(\omega),
\quad t\geq T_{\mathfrak B}(\omega),
\]
for any $\textbf{v}(0)\in\mathfrak{B}(\theta_{-t}\omega)$
 and $\textbf{u}(0)\in\mathcal{A}_0$,
where $T_{\mathfrak B}$ denotes the random variable defined in  Lemma \ref{lemma4.3}.
\end{lemma}

\begin{proof}

Taking the inner product of the system   \eqref{4.8} with $A^2w$ in $H$,  by integration by parts, we can deduce that
\begin{align}\label{4.17}
 &\frac12\frac{\d}{\d t}\|Aw\|^2- \nu \int_{\mathbb{T}^2}\partial_{yy}w_1\Delta^2 w_1dx-\nu \int_{\mathbb{T}^2}\partial_{xx}w_2\Delta^2 w_2dx\nonumber \\
&= \left(-\nu A_1hy_\delta(\theta_t\omega) +hy_\delta(\theta_t\omega),A^2w \right) \nonumber \\
&\quad -\left(B(\textbf{v}+hy_\delta(\theta_t\omega)) -B(\textbf{u}),A^2w \right)\nonumber\\
&=  \left(-\nu A_1hy_\delta(\theta_t\omega) +hy_\delta(\theta_t\omega),A^2w \right)\nonumber\\
&\quad -\left(B(\textbf{v}+hy_\delta(\theta_t\omega),w+hy_\delta(\theta_t\omega)),A^2w \right )\nonumber\\
&\quad
-\left(B(w+hy_\delta(\theta_t\omega),\textbf{u}),A^2w \right )\nonumber\\
&=: I_{4}(t)+I_{5}(t)+I_{6}(t).
\end{align}
By direct computations and $\nabla\cdot w=0$ and \eqref{3.9}, it also yields
\begin{align}
&-\nu \int_{\mathbb{T}^2}\partial_{yy}\Delta w_1\Delta w_1dx-\nu \int_{\mathbb{T}^2}\partial_{xx}\Delta w_2\Delta w_2dx\nonumber\\
&= \nu \int_{\mathbb{T}^2}(\partial_{y}\Delta w_1)^2dx
+\nu \int_{\mathbb{T}^2}(\partial_{x}\Delta w_2)^2dx
\nonumber\\
&\geq\frac12\nu\|A^{\frac32} w\|^2.
\end{align}

In order to get \eqref{sep6.3}, we only to estimate each terms on the right hand side of \eqref{4.17}.
By using H\"{o}lder's inequality, Young's inequality and $h\in H^3$,  we can deduce that
\begin{align}
I_{4}(t)&\leq  \nu \|A^{\frac32}hy_\delta(\theta_t\omega)\|\|A^{\frac32}w\|+ \|A^{\frac12}hy_\delta(\theta_t\omega)\|\|A^{\frac32}w\|\nonumber\\
&\leq \frac{ \nu }{16}\|A^{\frac32}w\|^2+C  |y_\delta(\theta_t\omega)|^2. \label{4.18}
\end{align}
By using H\"{o}lder's inequality and the Gagliardo-Nirenberg inequality, it also yields
\begin{align}
I_{5}(t)&\leq\|A^\frac12\textbf{u}_h\|_{L^4}\|\nabla(w+hy_\delta(\theta_t\omega))\|_{L^4}\|A^\frac32w\|
+\|\textbf{u}_h\|_{L^4}\|A(w+hy_\delta(\theta_t\omega))\|_{L^4}\|A^\frac32w\|
\nonumber\\
&\leq C\|A^\frac12\textbf{u}_h\|^\frac12\|A\textbf{u}_h\|^\frac12
\|\nabla(w+hy_\delta(\theta_t\omega))\|^\frac12\|A(w+hy_\delta(\theta_t\omega))\|^\frac12
\|A^\frac32w\|\nonumber\\
& \quad +C\|\textbf{u}_h\|^\frac12\|A^{\frac12}\textbf{u}_h\|^\frac12
\|A(w+hy_\delta(\theta_t\omega))\|^\frac12\|A^{\frac32}(w+hy_\delta(\theta_t\omega))\|^\frac12\|A^\frac32w\|\nonumber\\
&\leq  C\|A^\frac12\textbf{u}_h\|^\frac12\|A\textbf{u}_h\|^\frac12
\left(\|A^{\frac12}w\|^\frac12+ |y_\delta(\theta_t\omega)|^\frac12 \right) \left (\|Aw\|^\frac12+ |y_\delta(\theta_t\omega) |^\frac12 \right)
\|A^\frac32w\|\nonumber\\
&\quad +C\|A^{\frac12}\textbf{u}_h\| \left (\|Aw\|^\frac12+ |y_\delta(\theta_t\omega)|^\frac12\right)\left(\|A^{\frac32}w\|^\frac12
+ | y_\delta(\theta_t\omega)|^\frac12 \right)\|A^\frac32w\|\nonumber\\
&=: J_1(t) +J_2(t) . \nonumber
\end{align}

Since $\mathcal{A}_0$ is bounded in $H^2$,  $\|A\textbf{u}(t) \|^2+\|A^{1/2}\textbf{u}(t) \|^2\leq C$ for each $t\geq 0$, hence we can deduce that
\begin{align}\label{4.15}
\|A\textbf{u}_h\|^2&\leq C\|Aw\|^2+C\|A\textbf{u}\|^2
 +  C\|Ahy_\delta(\theta_t\omega)\|^2 \nonumber \\
&
\leq C\|Aw\|^2+  C  +C \|Ahy_\delta(\theta_t\omega)\|^2
,
\end{align}
and similarly, we also have
\begin{align}  \label{mar5.3}
  \|A^{\frac 12}w\|^2
   &\leq  C\|A^{\frac 12} \textbf{v}\|^2 +C
   \nonumber  \\
   &\leq C\|A^{\frac 12} \textbf{u}_h\|^2 + C\|A^{\frac 12} hy_\delta(\theta_t\omega) \|^2 +C .
\end{align}
Hence,  for  the term $J_1(t)$, we introduce the following inequalities
\begin{align*}
  &C\|A^\frac12\textbf{u}_h\|^\frac12\|A\textbf{u}_h\|^\frac12\|A^\frac12w\|^\frac12
\|Aw\|^\frac12\|A^\frac32w\|
& \nonumber\\
&\quad  \leq    \frac \nu {64}\|A^\frac32w\|^2  + C \|A^\frac12\textbf{u}_h\|  \|A \textbf{u}_h\|   \|A^{\frac 12} w\|  \|Aw\|
\nonumber
\\
&\quad  \leq   \frac \nu {64}\|A^\frac32w\|^2
+ C \|A^\frac12\textbf{u}_h\|^2  \|A \textbf{u}_h\|^2+      \|A^{\frac 12} w\|^2 \|Aw\|^2
\nonumber
\\
&\quad  \leq    \frac \nu {64} \|A^\frac32w\|^2
+ C\left( \|A^\frac12\textbf{u}_h\|^2 +  \|A^\frac 12 w   \|^2  \right) \|A w\|^ 2
+ C  \|A^\frac12\textbf{u}_h\|^2   \left(  1+ |y_\delta(\theta_t\omega) |^2  \right) ,
\end{align*}
\begin{align*}
  &C\|A^\frac12\textbf{u}_h\|^\frac12\|A\textbf{u}_h\|^\frac12\|A^\frac12w\|^\frac12
 |y_\delta(\theta_t\omega)|^\frac12\|A^\frac32w\|
& \nonumber\\
&\quad  \leq   \frac \nu {64}\|A^\frac32w\|^2  + C \|A^\frac12\textbf{u}_h\|  \|A \textbf{u}_h\|   \|A^{\frac 12} w\| |y_\delta(\theta_t\omega) |
\nonumber
\\
& \quad  \leq    \frac \nu {64}\|A^\frac32w\|^2
+ C \|A^\frac12\textbf{u}_h\|^2  \|A \textbf{u}_h\|^2+   C   \|A^{\frac 12} w\|^2 |y_\delta(\theta_t\omega) |^2
\nonumber
\\
&\quad \leq      \frac \nu {64}\|A^\frac32w\|^2
+ C  \|A^\frac12\textbf{u}_h\|^2  \|A w\|^2  + C \big(  \|A^\frac12\textbf{u}_h\|^2 +|y_\delta(\theta_t\omega) |^2 +1\big) \big(  1+|y_\delta(\theta_t\omega) |^2  \big)  ,
\end{align*}

\begin{align*}
  &C\|A^\frac12\textbf{u}_h\|^\frac12\|A\textbf{u}_h\|^\frac12 |y_\delta(\theta_t\omega)|^\frac12
\|Aw\|^\frac12\|A^\frac32w\|
& \nonumber\\
&\quad  \leq    \frac \nu {64} \|A^\frac32w\|^2  + C \|A^\frac12\textbf{u}_h\|  \|A \textbf{u}_h\|  |y_\delta(\theta_t\omega) |  \|Aw\|
\nonumber
\\
&\quad \leq     \frac \nu {64}\|A^\frac32w\|^2
+ C \|A^\frac12\textbf{u}_h\|^2  \|A \textbf{u}_h\|^2+    C| y_\delta(\theta_t\omega)|^2 \|Aw\|^2
\nonumber
\\
&\quad \leq     \frac \nu {64}\|A^\frac32w\|^2
+ C\left( \|A^\frac12\textbf{u}_h\|^2 +  |y_\delta(\theta_t\omega)|^2  \right) \|A w\|^2
+ C  \|A^\frac12\textbf{u}_h\|^2   \left(  1+ |y_\delta(\theta_t\omega) |^2  \right) ,
\end{align*}
and
\begin{align*}
  &C\|A^\frac12\textbf{u}_h\|^\frac12\|A\textbf{u}_h\|^\frac12 | y_\delta(\theta_t\omega)|  \|A^\frac32w\|
& \nonumber\\
&\quad \leq     \frac \nu {64}\|A^\frac32w\|^2  + C \|A^\frac12\textbf{u}_h\| \|A \textbf{u}_h\|   |y_\delta(\theta_t\omega)|^2
\nonumber
\\
&\quad \leq    \frac \nu {64}\|A^\frac32w\|^2
+ C \|A^\frac12\textbf{u}_h\|^2  \|A \textbf{u}_h\|^2   + C|y_\delta(\theta_t\omega) |^4
\nonumber
\\
&\quad \leq   \frac \nu {64}\|A^\frac32w\|^2
+ C  \|A^\frac12\textbf{u}_h\|^2     \left(  \|A w\|^2  +1+|y_\delta(\theta_t\omega) |^2  \right)
+ C|y_\delta(\theta_t\omega) |^4 ,
\end{align*}
it follows that
\begin{align}
J_1(t)
&\leq \frac  \nu {16}\|A^\frac32w\|^2
+C\left(\|A^\frac12\textbf{u}_h\|^2+\|A^\frac12w\|^2+|y_\delta(\theta_t\omega)|^2 \right)
\|Aw\|^2\nonumber\\
&\quad +C\left(1+\|A^\frac12\textbf{u}_h\|^4 \right)+C   |y_\delta(\theta_t\omega)|^4 .
\nonumber
\end{align}

For   the term $J_2(t)$, applying \eqref{4.15} and the Gagliardo-Nirenberg inequality and the Young inequality we can deduce that
\begin{align}
J_2(t)
&\leq C\|A^{\frac12}\textbf{u}_h\|\|Aw\|^\frac12\|A^{\frac32}w\|^\frac 32 +C\|A^{\frac12}\textbf{u}_h\|\|Aw\|^\frac12|y_\delta(\theta_t\omega)|^\frac12\|A^\frac32w\|
 \nonumber\\
&\quad+C\|A^{\frac12}\textbf{u}_h\| |y_\delta(\theta_t\omega)|^\frac12\|A^{\frac32}w\|^\frac 32
+C\|A^{\frac12}\textbf{u}_h\| |y_\delta(\theta_t\omega)|  \|A^\frac32w\|\nonumber\\
&\leq \frac \nu {16}\|A^\frac32w\|^2+C\left (1+\|A^\frac12\textbf{u}_h\|^4\right )\|Aw\|^2+C\left (1+\|A^\frac12\textbf{u}_h\|^8 \right)  +C |y_\delta(\theta_t\omega)|^4, \nonumber
\end{align}
and then  for $I_{5} (t)=J_1(t)+ J_2(t)$ we can get the following estimate
\begin{align}
I_{5}(t) &\leq  \frac \nu 8\|A^\frac32w\|^2
+C\left(1+ \|A^\frac12\textbf{u}_h\|^4+\|A^\frac12w\|^2+ |y_\delta(\theta_t\omega)|^2 \right)
\|Aw\|^2 \nonumber  \\
&\quad +C\left(1+\|A^\frac12\textbf{u}_h\|^8 \right)
+C  |y_\delta(\theta_t\omega)|^4. \label{sep6.2}
\end{align}


For the term $I_{6}(t)$, applying the similar method, we also can deduce that
\begin{align}\label{4.16}
I_{6}(t)&\leq\|A^{\frac12}(w+hy_\delta(\theta_t\omega))\|_{L^4}\|\nabla \textbf{u}\|_{L^4}\|A^\frac32w\|+\|w+hy_\delta(\theta_t\omega)\|_{L^\infty}\|A\textbf{u}\|\|A^\frac32w\|\nonumber\\
&\leq C\|A^{\frac12}(w+hy_\delta(\theta_t\omega))\|^\frac12\|A(w+hy_\delta(\theta_t\omega))\|^\frac12\|A^\frac12\textbf{u}\|^\frac12
\|A\textbf{u}\|^\frac12\|A^\frac32w\|\nonumber\\
& \quad +C\|A^\frac12(w+hy_\delta(\theta_t\omega))\|^\frac12\|A(w+hy_\delta(\theta_t\omega))\|^\frac12\|A\textbf{u}\|\|A^\frac32w\|\nonumber\\
&\leq C\|A^{\frac12}(w+hy_\delta(\theta_t\omega))\|^\frac12\|A(w+hy_\delta(\theta_t\omega))\|^\frac12\|A^\frac32w\|\nonumber\\
&\leq \frac \nu {16}\|A^\frac32w\|^2+C\left(\|A^{\frac12}w\|+\|A^{\frac12}hy_\delta(\theta_t\omega)\| \right) \big (\|Aw\|
+\|Ahy_\delta(\theta_t\omega)\| \big )\nonumber\\
&\leq \frac \nu {16}\|A^\frac32w\|^2+C\|Aw\|^2 +C |y_\delta(\theta_t\omega)|^2.
\end{align}
Inserting \eqref{4.18}, \eqref{sep6.2} and \eqref{4.16} into \eqref{4.17}, by \eqref{mar5.3}  we can get that
\begin{align}
 \frac{\d}{\d t}\|Aw\|^2+  \frac{\nu}{2}  \|A^{\frac32}w\|^2
&\leq C\left (1+\|A^\frac12\textbf{u}_h\|^4+\|A^\frac12w\|^2+|y_\delta(\theta_t\omega)|^2 \right )
\|Aw\|^2\nonumber\\
&\quad +C \left (1+\|A^\frac12\textbf{u}_h\|^8+|y_\delta(\theta_t\omega)|^4 \right) \nonumber \\
 &\leq C\left (1+\|A^\frac12 \textbf{v}\|^4 +|y_\delta(\theta_t\omega)|^4 \right )
\|Aw\|^2  \nonumber \\
&\quad +C \left (1+\|A^\frac12 \textbf{v}\|^8 +|y_\delta(\theta_t\omega)|^8 \right) .  \label{sep6.3}
\end{align}

 For any $t >1$ and $s\in[t-1,t]$, applying Gronwall's lemma to \eqref{sep6.3}   it is easy to get that
\begin{align*}
\|Aw(t)\|^2&\leq e^{C\int_s^t( 1+\|A^\frac12 \textbf{v}(\eta)\|^4 +|y_\delta(\theta_\eta\omega)|^4)\, \d \eta}\|Aw(s)\|^2 \nonumber\\
&\quad +\int_s^t Ce^{C\int_\tau^t  (1+\|A^\frac12 \textbf{v}\|^4+|y_\delta(\theta_\eta\omega)|^4)\, \d \eta}
 \Big (1+\|A^\frac12 \textbf{v}\|^8  +|y_\delta(\theta_\tau\omega)|^8\Big ) \, \d \tau\nonumber\\
&\leq Ce^{C\int_{t-1}^t(1+\|A^\frac12 \textbf{v}\|^8 +|y_\delta(\theta_\eta\omega)|^8)\, \d \eta} \nonumber \\
&\quad \times \left [\|Aw(s)\|^2+\int_{t-1}^t
\left(1+\|A^\frac12 \textbf{v}\|^8  +|y_\delta(\theta_\tau\omega)|^8\right)\d \tau \right]\\
& \leq Ce^{C\int_{t-1}^t(1+\|A^\frac12 \textbf{v}\|^8 +|y_\delta(\theta_\eta\omega)|^8)\, \d \eta} \Big(\|Aw(s)\|^2+ 1 \Big)
, \nonumber
\end{align*}
where in the last inequality we have used the basic relation $x\leq e^x$ for $x\geq 0$.
Integrating w.r.t. $s$ over   $  [t-1,t] $, we can get that
 \begin{align*}
\|Aw(t)\|^2
&\leq Ce^{C\int_{t-1}^t  (1+\|A^\frac12 \textbf{v}(\eta)\|^8+|y_\delta(\theta_\eta\omega)|^8)\, \d \eta}
 \left (\int_{t-1}^t\|Aw(s)\|^2\, \d s + 1\right). \nonumber
\end{align*}
Then we replace $\omega$ with $\theta_{-t}\omega$ to deduce
\begin{align}\label{4.22}
\|Aw(t,\theta_{-t}\omega,w(0))\|^2
&\leq Ce^{C\int_{t-1}^t \|A^\frac12 \textbf{v}(\eta,\theta_{-t}\omega, \textbf{v}(0))\|^8 \d \eta + C\int^0_{-1} |y_\delta(\theta_{\eta}\omega) |^8
\, \d \eta}\nonumber\\
&\quad \times \left(\int_{t-1}^t\|Aw(s,\theta_{-t}\omega,w(0))\|^2\, \d s +1  \right) .
\end{align}
By \eqref{mar8.3}, it can get that
\begin{align} \label{4.23}
  \int_{t-1}^t  \|A^\frac12 \textbf{v}(s,\theta_{-t}\omega,\textbf{v}(0))\|^8\, \d s \leq  |\zeta_3(\omega)|^4   ,  \quad t\geq T_{\mathfrak B}(\omega).
\end{align}
Moreover, since $\textbf{u}$ is a trajectory within the global attractor $\mathcal A_0$ which is   bounded  in $H^2$,  by \eqref{mar8.3}  again we can deduce that
\begin{align}
\int_{t-1}^t\|Aw(s,\theta_{-t}\omega,w(0))\|^2\, \d s
&\leq2\int_{t-1}^t\|A\textbf{v}(s,\theta_{-t}\omega,\textbf{v}(0))\|^2\, \d s+2\int_{t-1}^t\|A\textbf{u}\|^2\, \d s\nonumber\\
&\leq  2\int_{t-1}^t \|A\textbf{v}(s,\theta_{-t}\omega,\textbf{v}(0))\|^2\, \d s  +  2\|\mathcal A \|_{H^2} \nonumber\\
&\leq \frac 4  \nu \, \zeta_3(\omega) +C ,  \quad t\geq T_{\mathfrak B}(\omega).  \label{mar8.7}
\end{align}
Hence, inserting \eqref{4.23} and \eqref{mar8.7} into \eqref{4.22} we can deduce that
\begin{align}
\|Aw(t,\theta_{-t}\omega,w(0))\|^2
&\leq C e^{C|\zeta_3(\omega) |^4 + C\int^0_{-1} |y_\delta(\theta_{\eta}\omega) |^8
\, \d \eta} \left(   \zeta_3(\omega) +1
\right)  \nonumber \\
&=:\zeta_5(\omega),  \quad t\geq T_{\mathfrak B}(\omega).  \nonumber
\end{align}
 This completes the proof of Lemma \ref{lem:H2bd}.
\end{proof}

\end{document}
